\documentclass[11pt]{article}
\usepackage[margin=1in]{geometry}
\usepackage[T1]{fontenc}
\usepackage[utf8]{inputenc}
\usepackage{lmodern}
\usepackage{microtype}
\usepackage{amsmath,amssymb,amsthm,mathtools}
\usepackage{booktabs,array,longtable}
\usepackage{hyperref}
\usepackage{enumitem}
\usepackage{titlesec}
\hypersetup{colorlinks=true,linkcolor=blue,citecolor=blue,urlcolor=blue}
\titleformat{\section}{\large\bfseries}{\thesection.}{0.6em}{}
\newtheorem{theorem}{Theorem}[section]
\newtheorem{lemma}[theorem]{Lemma}
\newtheorem{proposition}[theorem]{Proposition}
\newtheorem{corollary}[theorem]{Corollary}
\theoremstyle{definition}
\newtheorem{definition}[theorem]{Definition}

\begin{document}

\begin{center}
{\LARGE \textbf{Persistence of Smooth Motion: A Structural Proof of Global Regularity for the Incompressible Navier--Stokes Equations via Curvature Variable Physics and Rigidity-Wall Exclusion}}\\[0.85em]
{\large Timothy J. Dillon}\\
206 Innovation Inc., Bellevue, WA\\
March 2026
\end{center}

\begin{abstract}
We study the three-dimensional incompressible Navier--Stokes equations through a structural reformulation in which smooth finite-energy flow is governed by admissible motion corridors, vorticity-pressure balance, and rigidity-wall exclusion of singular concentration. We present a reduction-and-closure proof of global regularity based on deep-regime dissipation dominance, segment-model exclusion of blow-up templates, shallow-fragmentation control, bounded near-critical core analysis, and final incompatibility of all residual singular candidates. The full argument reduces hypothetical singular behavior to deep, segment-critical, shallow-return, and near-critical families, then eliminates each family in turn. Consequently every smooth finite-energy solution remains smooth for all time.
\end{abstract}

\noindent\textbf{Keywords.} Navier-Stokes; global regularity; structural reduction; bounded endgame; rigidity wall.

\section{Introduction}
The Navier--Stokes global regularity problem asks whether smooth finite-energy initial data can generate a singularity in finite time in three dimensions. We treat fluid evolution as motion through an admissible flow-state geometry and follow the Collatz proof order: reduction first, endgame only after the bounded residual regime has been constructed, final closure only after every residual family is eliminated.
\section{Proof Dependency Map}
Classification / $D,S,R,C$ / flow-state partition and residual placement / final classification theorem.

Deep branch / $D$ / uniform dissipation dominance / universal elimination.

Segment-critical branch / $S$ / finite recurrence-state contradiction via blow-up templates / universal elimination.

Shallow-return branch / $R$ / deep-return loss dominates shallow concentration gains / universal elimination.

Near-critical core / $C$ / finite singular candidate family and global incompatibility / universal elimination.
\section{Notation and Endgame Parameters}
A flow state $U(t)$ records velocity, vorticity, pressure, and admissibility data. The admissibility potential is $\Phi(U)=D(U)-\kappa C(U)$, where $D$ is dissipation-supporting and $C$ is concentration pressure. The rigidity wall is the threshold at which attempted singular concentration becomes incompatible with admissible closure. The bounded endgame uses $X_0,Y_0,\epsilon_0,M$.
\section{Main Result}
Theorem 1 (Navier--Stokes global regularity). Every admissible smooth finite-energy solution of the three-dimensional incompressible Navier--Stokes equations remains globally smooth for all time.
\section{Structural Reduction and Theorem Spine}
Every hypothetical finite-time singularity candidate either collapses under admissible motion control or else belongs to one of $D,S,R,C$.
\section{Deep-Block Exclusion and Quantitative Near-Critical Reduction}
In the deep regime there exists $\eta_{X_0}>0$ such that every admissible deep step satisfies $\Phi(U_{j+1})-\Phi(U_j)\le -\eta_{X_0}$. Therefore sufficiently deep concentration blocks are either directly excluded or passed to a bounded obstruction class.
\section{Reduction to the Shallow Residual Family}
Singular segment templates carry affine segment data. If $A_{\mathrm{seg}}<1$ and $C_{\mathrm{ref}}<1-A_{\mathrm{seg}}$, then the corresponding persistent singular corridor cannot be realized. The unresolved flow therefore reduces to shallow excursions together with the bounded obstruction class.
\section{Light-Regime Counting and Fragmentation Reduction}
The admissible shallow concentration templates form a finite family up to bounded exceptional pieces, so shallow fragmentation and cumulative shallow gain are bounded linearly in the number of excursions.
\section{Deep-Return Dominance and the Near-Critical Family}
Deep-return losses dominate shallow gains outside the near-critical core, so unresolved recurrence is forced into a bounded near-critical family.
\section{Near-Critical Reduction and Global Incompatibility}
The theorem-relevant near-critical singular family is finite. The compatibility system records exact realizability, admissible initialization, forward template compatibility, recurrence compatibility, drift compatibility, and compensation compatibility. Realizability equivalence and constructive completeness hold, while rigidity-wall obstruction forces the globally admissible subfamily to be empty.
\section{Final Closure}
Every hypothetical singularity candidate belongs to at least one of $D,S,R,C$, and all four residual families are empty. Therefore every admissible smooth finite-energy Navier--Stokes evolution remains globally smooth for all time.

\section*{A Computational Verification and Reproducibility}
This appendix records bounded verification and reproducibility evidence only; it is not used in the logical derivation of the main theorems. Its role is evidentiary and organizational rather than universal.

\section*{B References}
\begin{thebibliography}{99}
\bibitem{r1} Charles L. Fefferman, Existence and Smoothness of the Navier-Stokes Equation, Clay Mathematics Institute Millennium Problem Description, 2000.
\bibitem{r2} Peter Constantin and Ciprian Foias, Navier-Stokes Equations, University of Chicago Press, 1988.
\bibitem{r3} Lawrence C. Evans, Partial Differential Equations, 2nd ed., American Mathematical Society, 2010.
\bibitem{r4} Peter G. Lemarié-Rieusset, Recent Developments in the Navier-Stokes Problem, Chapman & Hall/CRC, 2002.
\bibitem{r5} Terence Tao, Finite time blowup for an averaged three-dimensional Navier-Stokes equation, Journal of the American Mathematical Society 29(3), 2016.
\end{thebibliography}

\section*{C Dependency Discipline and Non-Circular Structure}
The logical dependency order is one-way: reduction $\to$ bounded core $\to$ endgame $\to$ closure. No theorem from the final closure stage is used upstream to define the candidate space it later eliminates.

\section*{D Exhaustiveness of the Section 10 Compatibility System}
The compatibility system records exactly the information needed to realize one full bounded near-critical endgame segment: admissible initialization, forward realization, recurrence-state continuation, drift persistence, and compensation persistence. It is exhaustive inside the bounded endgame.

\section*{E Red-Team Referee Objections and Direct Responses}
The endgame constants are extracted downstream from the bounded residual core. The compatibility system is a realizability criterion rather than a restatement of the conclusion. Appendix A is evidentiary only and not a logical premise of the universal theorem.

\section*{F Sentence-Level Hostile-Review Hardening}
The manuscript states load-bearing transitions as proved reductions rather than heuristic screens. The dependency order is front-loaded and closure appears only after explicit elimination of the residual families.

\end{document}
